\chapter{小参数中的计算机证明，实际证明了什么}\label{app:small}
\section{先列出缺口，避免把小例子当作全部证明}
五平均只需额外处理十六、十七、十八元；七平均需要二十二至二十六元。其他位置数已经由一般范围覆盖。

其中五平均的三个点与七平均的二十三至二十六元共用一条路线：假设所有安全候选都发生碰撞，把碰撞写成线性方程，再证明这些方程与零和、同余见证不相容。七平均二十二元另用矩阵群和一个带尺度标签的有限图。

因此这里有两种有限证书：一类穷尽线性方程\emph{模式}，另一类穷尽一个已经证明足够的有限\emph{状态空间}。两者都不对实际输入数值规定高度上限。

\section{一条碰撞等式，怎样变成线性方程}
设两个锚点为 $a,b$，令 $\Delta=a-b$。把每个轻值写成
\[x_i=b+\Delta t_i.
\]
这样 $b$ 对应参数零，$a$ 对应参数一。选择一个轻值子集 $S$，再加 $i$ 份 $a$、若干份 $b$ 补足 $p$ 项，均值为
\[b+\frac{\Delta}{p}\left(i+\sum_{j\in S}t_j\right).
\]
若它恰好碰撞到 $a$，就必须有 $\sum_{j\in S}t_j=p-i$。这是关于实际有理参数的严格线性方程，不只是模 $p$ 同余。

实际 $t_i$ 有无穷多个，但它们的余数 $z_i=t_i\pmod p$ 只有有限种。每个余数模式决定哪些候选和必须写哪些方程。于是可以对有限种模式逐一做精确消元，而消元得到的是每个系统的\emph{全部有理解}。

\section{保护集中的每个位置，为什么有一个私人素数}
先为危险支持各留一个位置，然后删掉所有多余的保护位置，得到一个“再删任何一项都不行”的集合 $T$。对其中某个位置 $v$，总有一个危险素数的例外支持只由它保护；否则 $v$ 还能删。这就是私人素数。

在这个素数处，其他保护位置与锚点 $b$ 同余，只有 $v$ 不同。因此，如果未保护部分相对 $b$ 的总和为零，整个配置的相对总和就只剩一个非零项，违背全局零和。这个矛盾只需要知道\emph{总和}，不必强迫每个轻值都与 $b$ 同余。
\InsertProof{lem:collision-small}

\section{如何读一行碰撞表}
五平均表的一行是 $c=6,\alpha=1,\beta=0$，余数列 $(3,4,4,4,4)$，唯一实际解为 $(-7,4,4,4,4)$。这些有理未知数一旦全部满足碰撞方程，就被迫取这五个值。加上锚点贡献，
\[D=5+1+(-7+4+4+4+4)=15.
\]
当位置数为十七或十八、存在保护位置时，私人素数整除 $\Delta$，未保护相对总和 $\Delta D$ 便在该素数处为零，正好触发上面的矛盾。

这不是验证了五个具体输入，而是证明一种可能的失败模式只能有这个解，然后把它排除。
\InsertProof{prop:collision-certificate}
\InsertTables{17_five_seven.tex}

\section{二十二元的3872个状态是什么}
控制模数为 $M=1936$。列记为 $(z,y)$，其中 $y$ 是模 $M$ 单位。先记录比值 $x=z/y\pmod M$，有1936种。

但实际可用的环 $R=\Z[1/35]$ 的单位模 $M$ 只覆盖全部单位的一半，记这半为 $U$，另一半为 $3U$。所以再加一位 $\epsilon$：$y\in U$ 时为零，$y\in3U$ 时为一。总共 $1936\times2=3872$ 个状态。

这是为何单纯遍历1936个方向不够。两个方向相同的列，可能一个处于允许单位倍数类，另一个不在；必须看真实宏能否把它们之间的尺度差也调整掉。
\InsertProof{lem:seven22}

\section{有限图的边，怎样提升成无限输入的操作}
矩阵 $(a,b;c,d)$ 作用于列后，第二坐标乘上 $cx+d$。于是 $x$ 按分式线性变换更新，而 $\epsilon$ 要加上 $cx+d$ 属于哪一类的标记。

证书从终端 $(0,0)$ 出发，把每条允许边都展开，直到没有新状态；结果恰好覆盖3872个。生成元已经有真实逆，所以这条父边树可以反向走回终端。至于实际整数有多大，最后由 $\Gamma(M,R)$ 的精确纤维引理处理。
\InsertProof{prop:seven-certificate}

“最大父边深度为八”只说有限标记图最多走八条这样的边，不能直接说原题只需八次平均。每条宏本身有若干原子操作，精确纤维输送也还需要平均词。有限图深度与原子词长度是两个不同计数。
